\documentclass{article}
\usepackage[fleqn]{amsmath}
\usepackage{bm}
\usepackage{amssymb}
\usepackage{graphicx}
\usepackage[title]{appendix}
\usepackage{gensymb}
\newcommand{\Rho}{\mathrm{P}}
\title{Marginal proba}
\author{Matthieu GENDRIN \\
	Orange Innovation
	}
\setlength{\parindent}{0cm}

\date{\today}
% Hint: \title{what ever}, \author{who care} and \date{when ever} could stand 
% before or after the \begin{document} command 
% BUT the \maketitle command MUST come AFTER the \begin{document} command! 
\begin{document}

\maketitle

\section{Marginal proba of a 3D gaussian}

\[\Sigma^{-1} = \begin{pmatrix}\ensuremath{\mathrm{S_x}} & a & b\\
a & \ensuremath{\mathrm{S_y}} & c\\
b & c & \ensuremath{\mathrm{S_z}}\end{pmatrix}\]

\begin{align}
    X^T \Sigma^{-1} X & = \frac{(z-m_z)^2}{\sigma_z^2} + \frac{(y-m_y)^2}{\sigma_y^2} + \frac{x^2}{\sigma_x^2} \nonumber \\
    \text{With:} & \nonumber \\
    \sigma_z^2 & = 1 / \ensuremath{\mathrm{S_x}} \nonumber \\
    m_z & = - \frac{(b x + c y)}{\ensuremath{\mathrm{S_z}}} \nonumber \\
    \sigma_y^2 & = \ensuremath{\mathrm{S_z}} / (\ensuremath{\mathrm{S_z}} * \ensuremath{\mathrm{S_y}} - c^2) \nonumber \\
    m_y & = \frac{\ensuremath{\mathrm{S_z}} (\frac{b c}{\ensuremath{\mathrm{S_z}}}-a) x}{\ensuremath{\mathrm{S_y}} \ensuremath{\mathrm{S_z}}-c^2} \nonumber \\
    \sigma_{x}^2 & = \frac{\ensuremath{\mathrm{S_z}}}{\ensuremath{\mathrm{S_x}} \ensuremath{\mathrm{S_z}} - b^2 - \frac{(b c - a \ensuremath{\mathrm{S_z}})^2}{\ensuremath{\mathrm{S_z}} \ensuremath{\mathrm{S_y}} - c^2}} \nonumber \\
\end{align}

I checked (with maxima), that $\sigma_x^2$ is the $[0,0]$ element of the $\Sigma$ matrix.

Since $m_z, \sigma_z$ don't depend on $z$, and $m_y, \sigma_z$ only depend on $x$. One can integrate through $z$ then $y$ the gaussian:
\begin{align}
    G(x,y,z) & = f \exp(-\frac{1}{2} (X^T \Sigma^{-1} X))
\end{align}
Thus, $\sigma_x$ is the variance of the marginal proba along $x$.

\begin{align}
    \Sigma_{0,0} & = \frac{\ensuremath{\mathrm{S_y}}\, \ensuremath{\mathrm{S_z}}-{{c}^{2}}}{\ensuremath{\mathrm{S_x}}\, ( \ensuremath{\mathrm{S_y}}\, \ensuremath{\mathrm{S_z}}-{{c}^{2}}) +a\, ( b c-\ensuremath{\mathrm{S_z}} a) +b\, ( a c-\ensuremath{\mathrm{S_y}} b) } \nonumber
\end{align}

\end{document}
